\documentclass[10pt,a4paper]{amsart}
\usepackage[margin=19mm]{geometry}
\usepackage{amsmath,amssymb,mathtools}
\usepackage[scr=rsfs,cal=euler]{mathalfa}
\usepackage{xfrac}
\usepackage[unicode,hidelinks,bookmarks=false]{hyperref}
\newcommand{\ie}{i.e.\ }
\newcommand{\cf}{cf.\ }
\newcommand{\resp}{respectively\ }
\newcommand{\Subsection}{\subsection}
\providecommand{\nobreakdash}{\nobreak\hbox{-}}
\setlength{\emergencystretch}{2em}
\multlinegap0pt
\usepackage{mdframed}
\usepackage{mdframed}
\usepackage{mdframed}
\usepackage{mdframed}
\usepackage{tcolorbox}
\usepackage{amssymb}
\usepackage{bm}
\usepackage{xcolor}
\usepackage{float}
\usepackage{algorithm}\usepackage{algpseudocode}
\usepackage{mdframed}
\usepackage[shortlabels]{enumitem}
\newtheorem{lemma}{Lemma}
\def\BR{{\mathbb{R}}}
\def\fX{{\mathcal{X}}}
\def\fY{{\mathcal{Y}}}
\def\vg{\mathbf{g}}
\def\vu{\mathbf{u}}
\def\vx{\mathbf{x}}
\def\vy{\mathbf{y}}
\title{Nonsmooth Nonconvex-Concave Minimax Optimization: Convergence Criteria and Algorithms}
\author{Jinyang Shi, Luo Luo}
\date{Source: arXiv:2604.21371v1 (CC BY 4.0)}
\begin{document}
\maketitle

\noindent\emph{Source and licence.} Nonsmooth Nonconvex-Concave Minimax Optimization: Convergence Criteria and Algorithms. Version arXiv:2604.21371v1 (CC BY 4.0), distributed by arXiv under the Creative Commons Attribution 4.0 International licence, http://creativecommons.org/licenses/by/4.0/, as recorded in the arXiv licence metadata for this exact version and shown on the abstract page https://arxiv.org/abs/2604.21371v1. Full licence notice: \texttt{licenses/arxiv-2604.21371-CC-BY-4.0.txt}.\par\smallskip

\noindent\textit{Editorial scope.} Counted unit: the article's lemma g\&G\_2 (source0.tex lines 1542--1555). The article's complete original proof of it is delivered with it (source0.tex lines 1557--1579, one \texttt{proof} environment of 23 lines). No mathematics is written, repaired or re-derived: the statement and the proof are the article's own lines, in the article's own notation, and no retained line is substituted or re-flowed.  No further source-local context is needed: every cross-reference of every retained line resolves inside the two delivered spans.  Nothing was omitted from any retained span, so the package carries no omission record.  No retained line contains a citation, so the wrapper carries no bibliography and no bibliography entry.  No retained line contains a figure environment, an image-inclusion command or drawing code, so no image is delivered and the package declares no asset dependency.  Editorial formatting only, in addition: an AMS \texttt{amsart} preamble that restates the article's own declarations and macros used by the retained lines, namely the environments \texttt{lemma} on separate counters, together with the standard packages those lines need.  The retained environments print in this document as Lemma 1 in order of appearance, whereas the article prints the same items with the numbering of its own full document.  The complete native source of the article as distributed in the arXiv e-print for this exact version is bundled unchanged as \texttt{sources/199021/source0.tex}.
\begin{lemma}\label{lemma: g&G_2}
For all $\vx\in\fX$, $\vy\in\fY$, $\vg_x\in\BR^{d_x}$, $\vg_y\in\BR^{d_y}$, and $\eta>0$, it holds
\[
\langle \vg_x, G_x(\vx,\vy, \vg_x, \eta) \rangle \geq \|G_x(\vx,\vy, \vg_x, \eta)\|^2
\quad\text{and}\quad
\langle \vg_y, G_y(\vx,\vy, \vg_y, \eta) \rangle \geq \|G_y(\vx,\vy, \vg_y, \eta)\|^2.
\]
Applying Cauchy--Schwarz inequality, we also achieve
\[
\|\vg_x\| \geq \|G_x(\vx,\vy, \vg_x, \eta)\|
\quad\text{and}\quad
\|\vg_y\| \geq \|G_y(\vx,\vy, \vg_y, \eta)\|.
\]
\end{lemma}

\begin{proof}
We denote $\vx^+ := {\Pi}_\fX\left(\vx - \eta \vg_x\right)$. 
The convexity and closedness of set $\fX$ implies for all $\vu \in \fX$, it holds
\begin{equation*}
    \left\langle \vx^+ - (\vx - \eta \vg_x), \vu - \vx^+ \right\rangle \geq 0
\end{equation*}
Dividing above equation by $\eta^2$ and taking $\vu = \vx$, we achieve
\begin{equation*}
    \frac{1}{\eta^2}\left\langle \vx^+ - (\vx - \eta \vg_x), \vx - \vx^+ \right\rangle \geq 0.
\end{equation*}
By rearranging above inequality, the definition of $G_x(\vx,\vy, \vg_x, \eta)$ leads to
\begin{equation*}
\left\langle \vg_x, G_x(\vx,\vy, \vg_x, \eta) \right\rangle  
= \left\langle \vg_x, \frac{\vx - \vx^+}{\eta} \right\rangle  
\geq \frac{1}{\eta^2}\left\langle \vx- \vx^+  , \vx - \vx^+\right\rangle  
= \|G_x(\vx,\vy, \vg_x, \eta)\|^2.
\end{equation*}
Following the above derivation, we also achieve
\[
\langle \vg_y, G_y(\vx,\vy, \vg_y, \eta) \rangle \geq \|G_y(\vx,\vy, \vg_y, \eta)\|^2.
\]
Applying Cauchy--Schwarz inequality on above two results, we finish the proof.
\end{proof}

\end{document}
